\documentclass[11pt]{article}
\usepackage[a4paper,margin=1in]{geometry}
\usepackage{fontspec}
\usepackage[boldfont=false]{xeCJK}
\setCJKmainfont{FandolSong-Regular.otf}
\usepackage{amsmath}
\usepackage{amssymb}
\usepackage{array}
\usepackage{booktabs}
\usepackage{graphicx}
\usepackage{adjustbox}
\usepackage{enumitem}
\setlist[itemize]{topsep=2pt,partopsep=0pt,parsep=0pt,itemsep=2pt,leftmargin=1.4em}
\setlist[enumerate]{topsep=2pt,partopsep=0pt,parsep=0pt,itemsep=2pt,leftmargin=1.6em}
\usepackage{parskip}
\usepackage{fvextra}
\fvset{breaklines=true,breakanywhere=true,fontsize=\small}
\setlength{\parindent}{0pt}

\begin{document}

\begin{center}
  {\LARGE\bfseries Enzymes}\\[0.35em]
  {\large Igcse Biology}
\end{center}
\vspace{0.6em}

\section*{What enzymes are}

A \textbf{catalyst} 催化剂 is a substance that speeds up the \textbf{rate} 速率 of a \textbf{chemical reaction} 化学反应 but is \textbf{not} used up or changed by the reaction. The same catalyst can be used again and again.

\textbf{Enzymes} 酶 are biological catalysts — catalysts made by living cells. They are \textbf{proteins} 蛋白质. Enzymes take part in all the reactions of \textbf{metabolism} 新陈代谢, so they control almost every chemical reaction in a living thing.

Enzymes are vital. Without them, the reactions that keep you alive would be far too slow. Enzymes raise the reaction rate enough to keep the organism alive.

\section*{How an enzyme works}

Each enzyme has a special pocket called its \textbf{active site} 活性位点. The \textbf{molecule} 分子 it works on is called the \textbf{substrate} 底物.

\begin{itemize}
\item The shape of the active site is \textbf{complementary} 互补 to the shape of the substrate — they fit together like a key in a lock.

\item The substrate slots into the active site. \textbf{(Supplement)} This makes an \textbf{enzyme-substrate complex} 酶底物复合物.

\item The reaction takes place, and the substrate is changed into one or more \textbf{products} 产物.

\item The products leave the active site, which is then free to be used again.

\end{itemize}

\subsection*{Why enzymes are specific (Supplement)}

Each enzyme is \textbf{specific} 专一 — it works on only one kind of substrate. This is because only that substrate has the right shape to fit the active site. A substrate of a different shape will not fit, just as the wrong key will not open a lock.

\section*{The effect of temperature and pH}

The activity of an enzyme depends on \textbf{temperature} 温度 and pH. You can investigate this by measuring how fast the reaction goes at different temperatures or pH values.

\subsection*{Temperature}

\begin{itemize}
\item As the temperature rises from cold, enzyme activity speeds up. The enzyme and substrate molecules have more \textbf{kinetic energy} 动能, so they move faster and meet more often. \textbf{(Supplement)} There are more useful \textbf{collisions} 碰撞 each second.

\item Activity is highest at the \textbf{optimum temperature} 最适温度 (about 37 °C in the human body).

\item Above the optimum, activity falls fast. The heat changes the shape of the active site, so the substrate no longer fits. The enzyme is \textbf{denatured} 变性. Denaturation is permanent — the enzyme cannot recover.

\end{itemize}

\subsection*{pH}

\begin{itemize}
\item Each enzyme works best at one particular pH.

\item If the pH is too high or too low, the shape of the active site changes, the substrate stops fitting, and the enzyme is denatured.

\end{itemize}

\section*{Exam tips}

\begin{itemize}
\item An enzyme is both a \textbf{protein} and a biological catalyst: it speeds up a reaction and is \textbf{not} used up.

\item Describe the shape: the active site is \textbf{complementary} to the substrate (the lock-and-key idea). Each enzyme is \textbf{specific} to one substrate.

\item Low temperature → slow (too little kinetic energy). Optimum → fastest. Too hot → \textbf{denatured} (the active site loses its shape, permanently).

\item The wrong pH also \textbf{denatures} the enzyme by changing the shape of the active site.

\item "Denatured" does \textbf{not} mean "killed" — enzymes are not alive. Say that the shape of the active site has changed so the substrate no longer fits.

\end{itemize}


\end{document}
